\def\writeANS{}
\begin{loigiaibth}{14}
  \begin {center} \begin {longtable}{|c|c|c|c|c|c|} \hline Phân tử & CTCT & Liên kết & Phân tử & CTCT & Liên kết \\ \hline $H_2$ & H-H & 1H-H & $N_2H_4$ &\makecell [c]{\rule {0pt}{-2pt}\\ \chemfig {H-N(-[:-90]H)-N(-[:-90]H)-H}\\\rule {0pt}{-2pt}} & 1N-N, 4N-H \\ \hline $Cl_2$ & Cl-Cl & 1Cl-Cl & $CH_4$ &\makecell [c]{\rule {0pt}{2pt}\\ \chemfig {H-C(-[:-90]H)(-[:90]H)-H} \\ \rule {0pt}{2pt}} & 4C-H \\ \hline $O_2$ & O=O & 1O=O & $C_2H_6$ & \makecell [c]{\rule {0pt}{2pt}\\ \chemfig {H-C(-[:-90]H)(-[:90]H)-C(-[:-90]H)(-[:90]H)-H} \\ \rule {0pt}{2pt}} & 1C-C, 6C-H \\ \hline $N_2$ & N$\equiv $N & 1N$\equiv $N & $C_2H_4$ &\makecell [c]{\rule {0pt}{2pt}\\ \chemfig {H-C(-[:-90]H)=C(-[:-90]H)-H}\\ \rule {0pt}{2pt}}& 1C=C, 4C-H \\ \hline $HCl$ & H-Cl & 1H-Cl & $C_2H_2$ & $H-C\equiv C-H$ & 1C$\equiv $C, 2C-H \\ \hline $H_2O$ & H-O-H & 2H-O & $C_3H_8$ &\makecell [c]{\rule {0pt}{2pt}\\ \chemfig {H-C(-[:-90]H)(-[:90]H)-C(-[:-90]H)(-[:90]H)-C(-[:-90]H)(-[:90]H)-H}\\ \rule {0pt}{2pt}}& 2C-C, 8C-H \\ \hline $NH_3$ & \makecell [c]{\rule {0pt}{2pt}\\ \chemfig {H-N(-[:-90]H)-H}\\ \rule {0pt}{2pt}} & 3N-H & $CO_2$ & O=C=O & 2C=O \\ \hline \end {longtable} \end {center}
\end{loigiaibth}
\def\writeANS{}
\begin{loigiaibth}{15}
  \begin {align*} \Delta _r H_{298}^\circ \text { (1)}&= E_b(H - H) + E_b(Cl - Cl) - 2E_b(H - Cl)\\ &= 436 + 243 - 2.431 = -183 \text { kJ}. \\ \Delta _r H_{298}^\circ \text { (2)} &= E_b(N \equiv N) + 3.E_b(H - H) - 2.3E_b(N - H)\\ & = 946 + 3.436 - 2.3.389 = -80 \text { kJ}. \\ \Delta _r H_{298}^\circ \text { (3)} &= E_b(N - N) + 4E_b(N - H) - E_b(N \equiv N) - 2E_b(H - H) \\ &= 163 + 4.389 - 946 - 2.436 = -99 \text { kJ}. \\ \Delta _r H_{298}^\circ \text { (4)} &= 4E_b(H - Cl) + E_b(O = O) - 2E_b(Cl - Cl) - 2.2E_b(H - O) \\ &= 4.431 + 498 - 2.243 - 4.464 = -120 \text { kJ}. \\ \Delta _r H_{298}^\circ \text { (5)} &= 4E_b(C - H) + E_b(Cl - Cl) - 3E_b(C - H) - E_b(C - Cl) - E_b(H - Cl) \\ &= 4.414 + 243 - 3.414 - 339 - 431 = -113 \text { kJ}. \\ \Delta _r H_{298}^\circ \text { (6)} &= 2E_b(C - C) + 2.6E_b(C - H) + 8E_b(O = O) - 4.2E_b(C = O) - 6.2E_b(H - O) \\ &= 2.347 + 2.6.414 + 8.498 - 4.2.799 - 6.2.464 = -2314 \text { kJ}. \end {align*}
\end{loigiaibth}
\def\writeANS{}
\begin{loigiaibth}{16}
  \begin {enumerate} \item Ta có: $N \equiv N + O = O \longrightarrow 2N + 2O \longrightarrow 2NO$ \\ Vậy: $\Delta _r H_{298}^\circ = E_b(N \equiv N) + E_b(O = O) - 2E_b(NO)$ \\ $= 945 + 494 - 2.607 = + 225 \text { (kJ)}$ \item Năng lượng cần cung cấp để phá vỡ liên kết trong phân tử nitrogen và oxygen là rất lớn, trong khi phản ứng trên lại thu nhiệt nên chỉ có thể xảy ra ở nhiệt độ cao hoặc khi có tia lửa điện. \end {enumerate}
\end{loigiaibth}
\def\writeANS{}
\begin{loigiaibth}{17}
  \begin {align*} 3O_2(g) \longrightarrow 2O_3(g) \quad \Delta _r H_{298}^\circ &= 3 \times 498 - 2 \times (498 + 204) = +90 \text { kJ} \\ 2O_3(g) \longrightarrow 3O_2(g) \quad \Delta _r H_{298}^\circ &= 2 \times (498 + 204) - 3 \times 498 = -90 \text { kJ} \end {align*} Dựa vào kết quả tính toán cho thấy quá trình: \\ $O_2 \longrightarrow 2O$ và $3O_2 \longrightarrow 2O_3$ có $\Delta _r H_{298}^\circ > 0 \Rightarrow $ Chứng tỏ $O_3$ không có khả năng tồn tại. \\ Quá trình: $2O_3 \longrightarrow 3O_2$ có $\Delta _r H_{298}^\circ < 0$, chứng tỏ khả năng tồn tại của $O_2$
\end{loigiaibth}
\def\writeANS{}
\begin{loigiaibth}{18}
  \begin {enumerate} \item Trong hợp chất $CH_3-C \equiv CH$ , số liên kết C-H: 4; C-C: 2; C$\equiv $C: 1. \item Biến thiên enthalpy của phản ứng: \[ \text {CH}_3-\text {C} \equiv \text {CH} \text {(g)} + \text {H}_2\text {(g)} \xrightarrow [$\text {t}^\circ , \text {Pd/PbCO}_3$] \text {CH}_3\text {-CH}=\text {CH}_2 \text {(g)} \] \begin {eqnarray*} \Delta _r H_{298}^\circ &=& E_b(C\equiv C) + E_b(C-C) + 4 \times E_b(C-H)\\ &\vphantom {=}& + E_b(H-H) - E_b(C=C) - E_b(C-C) - 6 \times E_b(C-H)\\ &=& 839 + 347 + 4 \times 413 + 432 - 614 - 347 - 6 \times 413 = -169 \text { kJ} \end {eqnarray*} \par \end {enumerate}
\end{loigiaibth}
\def\writeANS{}
\begin{loigiaibth}{19}
  Phản ứng hóa học: $CO$(g) + $Cl_2$(g) $\xrightarrow [$C\text {than hoạt tính}$]$ $COCl_2$(g)(Phosgene)\\ $\Rightarrow \Delta _fH_{298}^\circ = \sum E_b \text {(cđ)} - \sum E_b \text {(sp)} \Rightarrow \Delta _f H_{298}^\circ = 1075 + 243 - 2 \times 339 - 745 = -105 \text {kJ}$
\end{loigiaibth}
\def\writeANS{}
\begin{loigiaibth}{20}
  $\mathrm {H_3C-CH_2-OH}$ có 1 liên kết C-C, 5 liên kết C-H, 1 liên kết C-O và 1 liên kết O-H \\ $\mathrm {H_3C-O-CH_3}$ có 6 liên kết C-H và 2 liên kết C-O \\ Quá trình đã cho có biến thiên enthalpy chuẩn là: \\ $\Delta _rH_{298}^0 = E_b(\mathrm {H_3C-CH_2-OH}) – E_b(\mathrm {H_3C-O-CH_3} )$ \\ $= (347 + 5{,}414 + 360 + 464) – (6{,}414 + 2{,}360) = 37\,\mathrm {kJ}$ \\ $\Delta _rH_{298}^0 > 0$ chứng tỏ ở điều kiện chuẩn $\mathrm {H_3C-CH_2-OH}$ bền hơn $\mathrm {H_3C-O-CH_3}$
\end{loigiaibth}
\def\writeANS{}
\begin{loigiaibth}{21}
  \begin {enumerate} \item Hydrazine có công thức cấu tạo: $\mathrm {H_2N-NH_2}$. \\ Một phân tử hydrazine có 1 liên kết đơn $\mathrm {N-N}$ ($E_b = 160\,\mathrm {kJ/mol}$); 4 liên kết đơn $\mathrm {N-H}$ ($E_b = 391\,\mathrm {kJ/mol}$). \\ $\mathrm {N_2}$ có 1 liên kết ba $\mathrm {N\equiv N}$ ($E_b = 945\,\mathrm {kJ/mol}$), \\ $\mathrm {H_2}$ Có 1 liên kết đơn $\mathrm {H-H}$ ($E_b = 432\,\mathrm {kJ/mol}$). \\ Áp dụng công thức tính $\Delta _rH_{298}^0$ theo năng lượng liên kết: \\ $\Delta _rH_{298}^0 = E_b(\mathrm {N-N}) + 4.E_b(\mathrm {N-H}) - E_b(\mathrm {N\equiv N}) - 2.E_b(\mathrm {H-H})$ \\ $= 160 + 4 \times 391 - 945 - 2 \times 432 = - 85\,\mathrm {kJ}$ \item \begin {itemize} \item $\mathrm {N_2H_4}$ là chất lỏng ở điều kiện thường nên dễ bảo quản (nếu là chất khí cần nén ở áp suất cao gây nguy hiểm). \item Khối lượng riêng nhỏ nên nhẹ, phù hợp với nhiên liệu động cơ tên lửa (nếu nặng sẽ gây tốn năng lượng). \item $\Delta _rH_{298}^0 = - 85\,\mathrm {kJ}$ nên phản ứng có thể tự xảy ra mà không cần nguồn nhiệt ngoài. \item Giả sử 1 mol $\mathrm {N_2H_4}$ lỏng phản ứng (có thể tích khá nhỏ) sẽ sinh ra 3 mol khí có thể tích lớn hơn rất nhiều nên sẽ tạo được luồng khí đẩy tên lửa đi. \end {itemize} \end {enumerate}
\end{loigiaibth}
\def\writeANS{}
\begin{loigiaibth}{22}
  \begin {enumerate} \item a, $\mathrm {H_3C-CH_2-CH_2-CH_3 \longrightarrow CH_2=CH-CH=CH_2 + 2H_2}$ \\ $\Delta _rH_{298}^0 = 10E_b(\mathrm {C-H}) + 3E_b(\mathrm {C-C}) - 6E_b(\mathrm {C-H}) - 2E_b(\mathrm {C=C}) - E_b(\mathrm {C-C}) - 2E_b(\mathrm {H-H})$ \\ $= 4{,}414 + 2{,}347 - 2{,}611 - 2{,}436 = 256\,\mathrm {kJ}$ \item b, $6\mathrm {CH_4 \longrightarrow C_6H_6} + 9\mathrm {H_2}$ \\ $\Delta _rH_{298}^0 = 6.4\, E_b(\mathrm {C-H}) - 3E_b(\mathrm {C-C}) - E_b(\mathrm {C=C}) - 6E_b(\mathrm {C-H}) - 9.E_b(\mathrm {H-H})$ \\ $= 18{,}414 - 3{,}347 - 3{,}611 - 9{,}436 = 654\,\mathrm {kJ}$ \\ Các phản ứng này không thuận lợi về phương diện nhiệt \\ Phản ứng theo chiều ngược lại thuận lợi hơn về phương diện nhiệt: \\ \begin {tabular}{cccc} $\mathrm {CH_2=CH-CH=CH_2}$ + $\mathrm {2H_2}$& $\longrightarrow $& $\mathrm {H_3C-CH_2-CH_2-CH_3}$,& $\Delta _rH_{298}^\circ = -256\,\mathrm {kJ}$ \\ $\mathrm {C_6H_6}$ + $\mathrm {9H_2}$ &$\longrightarrow $& $\mathrm {6CH_4}$, &$\Delta _rH_{298}^\circ = -654\,\mathrm {kJ}$ \end {tabular} \end {enumerate}
\end{loigiaibth}
\def\writeANS{}
\begin{loigiaibth}{23}
  - Xét X là F: \\ $\mathrm {CH_4(g) + F_2(g) \longrightarrow CH_3F(g) + HF(g)}$ \\ $\Delta _rH_{298}^0 = 1 \times E_b (\mathrm {CH_4}) + 1 \times E_b (\mathrm {F_2}) - 1 \times E_b (\mathrm {HF}) - 1 \times E_b (\mathrm {CH_3F})$ \\ $\Delta _rH_{298}^0 = 1 \times 4E_{\mathrm {C-H}} + 1 \times E_{\mathrm {F-F}} - 1 \times E_{\mathrm {H-F}} - 1 \times (3E_{\mathrm {C-H}} + E_{\mathrm {C-F}})$ \\ $\Delta _rH_{298}^0 = 1 \times 4 \times 414 + 1 \times 159 - 1 \times 565 - 1 \times (3 \times 414 + 1 \times 485) = -477\,\mathrm {kJ}$ \\ - Xét X là Cl: \\ $\mathrm {CH_4(g) + Cl_2(g) \longrightarrow CH_3Cl(g) + HCl(g)}$ \\ $\Delta _rH_{298}^0 = 1 \times E_b (\mathrm {CH_4}) + 1 \times E_b (\mathrm {Cl_2}) - 1 \times E_b (\mathrm {HCl}) - E_b (\mathrm {CH_3Cl})$ \\ $\Delta _rH_{298}^0 = 1 \times 4E_{\mathrm {C-H}} + 1 \times E_{\mathrm {Cl-Cl}} - 1 \times E_{\mathrm {H-Cl}} - 1 \times (3E_{\mathrm {C-H}} + E_{\mathrm {C-Cl}})$ \\ $\Delta _rH_{298}^0 = 1 \times 4 \times 414 + 1 \times 243 - 1 \times 431 - 1 \times (3 \times 414 + 1 \times 339) = -113\,\mathrm {kJ}$ \\ - Xét X là Br: \\ $\mathrm {CH_4(g) + Br_2(g) \longrightarrow CH_3Br(g) + HBr(g)}$ \\ $\Delta _rH_{298}^0 = 1 \times E_b (\mathrm {CH_4}) + 1 \times E_b (\mathrm {Br_2}) - 1 \times E_b (\mathrm {HBr}) - E_b (\mathrm {CH_3Br})$ \\ $\Delta _rH_{298}^0 = 1 \times 4E_{\mathrm {C-H}} + 1 \times E_{\mathrm {Br-Br}} - 1 \times E_{\mathrm {H-Br}} - 1 \times (3E_{\mathrm {C-H}} + E_{\mathrm {C-Br}})$ \\ $\Delta _rH_{298}^0 = 1 \times 4 \times 414 + 1 \times 193 - 1 \times 364 - 1 \times (3 \times 414 + 1 \times 276) = -33\,\mathrm {kJ}$ \\ - Xét X là I: \\ $\mathrm {CH_4(g) + I_2(g) \longrightarrow CH_3I(g) + HI(g)}$ \\ $\Delta _rH_{298}^0 = 1 \times E_b (\mathrm {CH_4}) + 1 \times E_b (\mathrm {I_2}) - 1 \times E_b (\mathrm {HI}) - E_b (\mathrm {CH_3I})$ \\ $\Delta _rH_{298}^0 = 1 \times 4E_{\mathrm {C-H}} + 1 \times E_{\mathrm {I-I}} - 1 \times E_{\mathrm {H-I}} - 1 \times (3E_{\mathrm {C-H}} + E_{\mathrm {C-I}})$ \\ $\Delta _rH_{298}^0 = 1 \times 4 \times 414 + 1 \times 151 - 1 \times 297 - 1 \times (3 \times 414 + 1 \times 240) = 28\,\mathrm {kJ}$ \\ $\Rightarrow $ Từ F đến I, tính phi kim giảm dần nên khả năng tham gia phản ứng giảm dần.
\end{loigiaibth}
\def\writeANS{}
\begin{loigiaibth}{24}
  \begin {enumerate} \item (a) Tính biến thiên enthalpy chuẩn của phản ứng dựa vào nhiệt tạo thành: \\ Phương trình hoá học của phản ứng: \\ $\mathrm {C_3H_8(g) + 5O_2(g) \xrightarrow {t^o} 3CO_2(g) + 4H_2O(g)}$ \\ $\Delta _fH_{298}^0$ : \hspace {0.5cm} -105 \hspace {1.5cm} 0 \hspace {1.5cm} -393{,}50 \hspace {1.5cm} -241{,}82 \\ $\Delta _rH_{298}^0 = 3 \times (-393{,}50) + 4 \times (-241{,}82) -105 = -2042{,}78\,\mathrm {kJ}$ \item (b)Tính biến thiên enthalpy chuẩn của phản ứng dựa vào năng lượng liên kết \\ Áp dụng công thức: $\Delta _rH_{298}^0 = \sum E_b (\mathrm {cd}) - \sum E_b (\mathrm {sp})$ \\ Phân tử $\mathrm {C_3H_8}$ số liên kết là: C-H: 8; C-C: 2. \\ $\Delta _rH_{298}^0 = 8 \times 413 + 2 \times 347 + 5 \times 498 - 3 \times 2 \times 745 - 4 \times 2 \times 467 = -1\,718\,\mathrm {kJ}$ \\ Biến thiên enthalpy chuẩn của phản ứng dựa vào nhiệt tạo thành có giá trị $-2042{,}78\,\mathrm {kJ}$ âm hơn so với biến thiên enthalpy chuẩn của phản ứng dựa vào năng lượng liên kết $-1718\,\mathrm {kJ}$. \end {enumerate}
\end{loigiaibth}
\def\writeANS{}
\begin{loigiaibth}{25}
  \begin {enumerate} \item (a) Xét phản ứng: $\mathrm {C_4H_{10}(g) + \dfrac {13}{2}O_2(g) \longrightarrow 4CO_2(g) + 5H_2O(g)}$ \item (b) $\Delta _rH_{298}^0 = 3.E_{\mathrm {C-C}} + 10.E_{\mathrm {C-H}} + 6{,}5.E_{\mathrm {O=O}} - 4.2.E_{\mathrm {C=O}} - 5.2.E_{\mathrm {O-H}}$ \\ $= 3.346 + 10.418 + 6{,}5.495 - 8.799 - 10.467 = - 2\,626{,}5\,\mathrm {(kJ)}$. \item (c) $Q = \dfrac {12.10^3.2626{,}5}{58} = 964\,163{,}4\,\mathrm {(kJ)}$ \\ Nhiệt cần đun 1 ấm nước: $2.10^3.4{,}2.(100 - 25) = 630\,000\,\mathrm {(J)} = 630\,\mathrm {(kJ)}$. \\ Số ấm nước: $\dfrac {964\,163{,}4.60\%}{630} = 918\,\mathrm {(ấm\,nước)}$. \end {enumerate}
\end{loigiaibth}
